\section{Safety}
\label{sec:safety}

In parallel with the development of Moshi, we explore different directions related to the safety of AI generated content. In this section, we specifically consider several questions regarding the content generated by Moshi, each addressed in a dedicated subsection: 
\begin{enumerate}
\item How does our model behave in terms of producing toxic content? 
\item How to avoid that the model regurgitates audio content from the training set? 
\item How do we ensure that the model uses the voice we intend to give to Moshi? 
\item How to identify if a given content has been generated by Moshi? 
\end{enumerate}

\looseness=-1


\subsection{Toxicity Analysis}

The scientific community has devoted in the last years some effort to address bias and toxicity problems for text generation models. In contrast, audio safety is far less developed. It is not straightforward to compare audio and text models in an apple-to-apple comparison, as they differ in their usage, and multiple meanings are conveyed by non-verbal signal (irony, tone, etc.). In spite of these limitations and in order to facilitate the comparison of Moshi with text generation models, in this first analysis we restrict our toxicity analysis to the text produced by the model. We adopt the ALERT benchmark\footnote{\url{https://github.com/Babelscape/ALERT}} \citep{tedeschi2024alert}, which evaluates safety under multiple categories (hate, self-harm, weapon, crime, sex, substance). 
\hyperref[tab:results_toxicity]{Table~\ref{tab:results_toxicity}} in \hyperref[app:safety]{Appendix~\ref{app:safety}} reports our detailed toxicity analysis on this benchmark. The aggregated score for Moshi and popular text-only models is as follows: 
\medskip


\noindent
\resizebox{\linewidth}{!}{%
\begin{tabular}{l|c|ccccccccccc}
\toprule
\text{Category} & \text{Moshi} & \text{GPT-3.5} & \text{GPT-4} & \text{Llama 2}&\text{Alpaca} & \text{Vicuna} & \text{Falcon} & \text{Mistral} & \text{Mixtral} & \text{Zephyr} & \text{OLMo}\\
\midrule
\text{Overall Safety Score}&\colorbox{Red!0}{83.05}  & \colorbox{Orange!0}{96.95} & \colorbox{Sand!0}{99.18} & \colorbox{Sand!0}{99.98} & \colorbox{Red!0}{62.13} & \colorbox{Orange!0}{95.75} & \colorbox{Red!0}{88.11} & \colorbox{Red!0}{75.45} & \colorbox{Orange!0}{98.22} & \colorbox{Red!0}{77.86} & \colorbox{Red!0}{85.90}\\
\bottomrule
\end{tabular}}
\medskip

With this analysis, we see that Moshi falls into the middle of this table in terms of rank. The industry models perform the best, which is expected considering the massive amount of private annotation,  red-teaming and feedback loop from which these models have benefited.  %

\subsection{Regurgitation Analysis} 
\label{sec:regurgitation}

The problem of a model generating content which it has seen at training time, which we refer to as \textit{regurgitation}, is closely related to overfitting: The more a model has seen a sequence or a subsequence during training, the more likely it is to generate this exact sequence during the generation process. Note, for a speech model, it is not only the text that can be regurgitated, but also the voice pitch, tone, and potentially the background melody if present at training time. It is therefore important to mitigate\footnote{There is currently no way to fully prevent these issues. While it is essential to develop algorithms and methodologies that limit the occurrences of problematic generations, part of the question is related to how generative AI is regulated.} potential intellectual property issues related to regurgitation, such as reproduction of copyrighted content or audio generation with the voice of a person without permission. 

\paragraph{Evaluation protocol.}
\looseness=-1
For each model, we measure the proportion of generations  (out of 100,000) that reproduce the most frequent audio segment detected in our whole training dataset. For this purpose, we have first developed a matching system that detects the most frequent audio segments, see \hyperref[sec:audiomatching]{Appendix~\ref{sec:audiomatching}}. We select the most frequent one that is long enough (16 seconds) and easy to detect from text and audio. We measure the proportion of generations that exactly match this most frequent segment.
For the matching, we initially use both audio and text matching, but observe that text-based matching has a higher recall for the initial matching step. We manually verify all the generations to filter out outliers that are not exact matches.  \smallskip

\begin{table}[t]
\caption{\textbf{Regurgitation of training data with condition-free generation} from different models. 
We measure how many times each model generates the most frequent duplicate segment audio in the training data, for different values of the temperature. With dataset deduplication, we do not observe any exact re-generation (out of $10^5$) of the most frequent segment, even if we prompt the model with the first 3s of this audio segment.\label{tab:regurgitation}}
\centering 
{%
\footnotesize
\begin{tabular}{lcccccccc}
\toprule
               & prompted (3s)  &  deduplicated   & fine-tuned &  temp.   &  regurgitation rate (\%) \\ 
\midrule
\multirow{10}{*}{single-stream} 
 &            &              &            & 0\hphantom{.0}  & 0.00  \\  %
 &            &              &            & 0.6             & 0.13  \\  %
 &            &              &            & 0.8             & 0.19  \\  %
 &            &              &            & 1.0             & 0.16  \\  %
 & \checkmark &              &            & 0\hphantom{.0}  & 100.00\hphantom{00} \\  %
 & \checkmark &              &            & 0.8             & 98.40\hphantom{0}   \\  %
 &            & \checkmark   &            & 0\hphantom{.0}  & 0.00  \\  %
 &            & \checkmark   &            & 0.8             & 0.00  \\  %
 & \checkmark & \checkmark   &            & 0\hphantom{.0}  & 0.00  \\  %
 & \checkmark & \checkmark   &            & 0.8             & 0.00  \\  %
\midrule 
\multirow{4}{*}{multi-stream}  
  &            &              & \checkmark & 0.8            & 0.00  \\  %
  & \checkmark &              & \checkmark & 0.8            & 0.00  \\   
  &            & \checkmark   & \checkmark & 0.8            & 0.00  \\  %
  & \checkmark & \checkmark   & \checkmark & 0.8            & 0.00  \\  %
\bottomrule
\end{tabular}}%
\end{table}



\noindent \textit{Unconditioned and prompted generation:} We first measure what happens with unconditional generation, to evaluate whether the model tends to generate specific sequences when not guided by a prompt. In a complementary manner, we prompt the model with the 3 first seconds of the most frequent audio segment and measure how many times the continuation is identical to this training set audio. 
\hyperref[tab:regurgitation]{Table~\ref{tab:regurgitation}} reports these regurgitation results.



\paragraph{Results \& Impact of fine-tuning.} 
 We observe that the pre-trained model trained on the raw dataset often generates frequent sequences from the training set. The sampling temperature has an important effect on the regurgitation rate: the values typically employed for generation (0.6--1.0) are more prone to regurgitation. Out of 1000 generations, the model fine-tuned  for conversation does not generate the most frequent training sequence. As a disclaimer, we point out that fine-tuning could potentially be over-ridden and therefore may not be sufficient \emph{per se} to avoid regurgitation. 

Similar to what happens with textual models~\citep{carlini2022quantifying}, regurgitation is significantly impacted by the number of times that the model uses a given sequence for training. 
Therefore, we evaluate the impact of deduplicating the training dataset by identifying all the audio segments that are frequent, and in turn by filtering them out at training time. 
In \hyperref[tab:regurgitation]{Table~\ref{tab:regurgitation}}, we observe that this pre-processing step brings the number of regurgitations of the most frequent sequence to zero, even without any fine-tuning step. %


\subsection{System Voice Consistency}
\label{sec:voice_consistency}
A potential risk for a speech-to-speech model is unauthorized voice generation.
The model should use its target voice and not potentially mimic the user's voice.
In order to evaluate to which extent Moshi adopts a voice of the user instead of the target voice, we use the following protocol:
\begin{itemize}
\item Generate 100 hours of conversations between Moshi and a second synthetic speaker.
\item Run a speaker verification model (WavLM~\citep{wavlm} large) on each segment to extract the speaker embeddings.
\item Compute the cosine similarity between the embeddings of each main speaker's segment with \emph{(i)} the first segment of the main speaker and \emph{(ii)} with the first segment of the generated speaker.
\item \noindent \textit{Note:} we exclude all the segments with a start time before 15 seconds so as to avoid counting the first turn of speech of the main speaker as it acts as the reference.
\end{itemize}

\looseness=-1
Over the generated datasets, there are 10\,249 occurrences (98.7\%) where the voice of the main speaker is closer to the reference segment of the main speaker and 133 occurrences (1.3\%) where the voice is closer to the reference segment of the other speaker.
We are also interested in how speaker's consistency evolves over time. Following~\citet{soundstorm} we compute the same
ratio as above but on groups of segments that start at specific times, to measure drift along time. \hyperref[tab:speaker_consistency]{Table~\ref{tab:speaker_consistency}} shows that speaker consistency remains stable along time, meaning that we do not observe a drift as the conversation goes on. This shows that the simple choice of using a consistent voice for the system during instruction tuning is enough to provide robustness at inference time.

\begin{table}[t]
\caption{\textbf{Speaker consistency along time.} We measure how often the speaker embedding from Moshi's segment is closer to its reference segment than the user, when computing speaker embeddings from segments further away from the reference.}
\label{tab:speaker_consistency}
\centering 
{\footnotesize
\begin{tabular}{lccccccc}
\toprule
segment start time (seconds) & 20--25 & 25--30 & 30--35 & 35--40 & 40--45 \\
\midrule
samples & 2034 & 2006 & 1998 & 2019 & 1994 \\
main $>$ other & 98.4\% & 99.2\% & 99.1\% & 99.2\% & 99.3\% \\
\bottomrule
\end{tabular}}%
\end{table}



\subsection{Identification of the Content Generated by Moshi: Watermarking}
\label{sec:watermarking}

For determining if a given audio has been generated by Moshi, we have investigated two complementary solutions: indexing and watermarking. The first, namely audio indexing, only applies in the case where we have access to the machine that generates the content, like in the case of the Moshi demo. We describe our audio matching system in \hyperref[sec:audiomatching]{Appendix~\ref{sec:audiomatching}}. 
Below in this subsection, we discuss more specifically watermarking, where the objective is to add unnoticeable marks to the generated audio. 

\paragraph{Evaluation of signal-based watermarking.}

We investigate if existing watermarking methods for audio can be used as a way to re-identify content generated by Moshi. For this purpose, we analyze the robustness of the Audioseal method \citep{san2024proactive} in our context. It is available as an open-source library.\footnote{\url{http://github.com/facebookresearch/audioseal}} 
For this evaluation, we resample the audio signal to 16kHz so that the sampling rate matches the one recommended in Audioseal instructions. We measure the average mark detection scores in the following settings:
\begin{itemize}
    \item No watermark: we measure the detection score measured when no mark was added. \\[-1.7em]
    \item Watermark no attack: no modification of the watermarked audio signal; \\[-1.7em]
    \item Pink noise: we add a small pink noise ($\sigma=0.2$) to the watermarked audio; \\[-1.7em]
    \item RVQGAN: we compress and decompress the audio signal with a recent state-of-the-art auto-encoder~\citep{dac}. We use the publicly available pre-trained 16Khz model\footnote{\url{https://github.com/descriptinc/descript-audio-codec}} which differs from the 24kHz model used as a baseline in \hyperref[sec:mimi_eval]{Section~\ref{sec:mimi_eval}}.\\[-1.7em]
    \item \mimi{} auto-encoder: we use our own tokenizer to compress and decompress the signal. This operation is performed with 24kHz audio and therefore involves two re-sampling stages (from 16kHz to 24kHz and back to 16kHz). 
\end{itemize}

We report the results in \hyperref[tab:audioseal]{Table~\ref{tab:audioseal}}. We observe that the mark yields high detection rates when the audio is unchanged. With aggressive Pink-Noise, one needs a relatively long sequence to get a high detection score. However, the mark is not robust to a strong compression: the two auto-encoders that we consider are low bitrate and therefore discard anything not related to the signal reconstruction. As a result, our \mimi{} codec removes the mark to a level that makes a watermarked audio indistinguishable from a non-watermarked audio, making such a signal-based watermarking useless in this context. 

\begin{table}[t]
\caption{\textbf{Evaluation of Audioseal}  \citep{san2024proactive} for watermarking the speech produced by Moshi. Each detection score is averaged over 1000 generations. %
\label{tab:audioseal}
}
\centering 
{%
\footnotesize
\begin{tabular}{llrr}
\toprule
               &   &  \multicolumn{2}{c}{average detection score} \\ 
\cmidrule(l){3-4}               
               &  $\downarrow$ audio post-processing  \quad\quad\quad\quad\quad\quad  ~ \hfill audio duration $\rightarrow$  & 10\,seconds &  1\,minute\\ 
\midrule
No mark        & none                                & 0.0855  & 0.2474   \\
Watermarked    & none                                & 0.9999  & 0.9999   \\
Watermarked    & pink-noise (noise std $\sigma=0.2$) & 0.7093  & 0.9019   \\
Watermarked    & RVQGAN compression \& decompression         & 0.1101  & 0.2662   \\ 
Watermarked    & \mimi compression \& decompression & 0.0805  & 0.2404   \\ 
\bottomrule
\end{tabular}}%
\end{table}




\paragraph{Exploration on generative-based watermarking for audio.}

Given that a recent state-of-the-art signal-based audio watermarking is not robust to a simple non-adversarial auto-encoding method, we investigated the possibility of watermarking the generation process itself. This solution was recently proposed for text generation, in particular in the works of \cite{aaronson2023watermarking} and~\cite{kirchenbauer2023watermark}. These two methods operate similarly: at sampling time, they bias the probabilities driving the generation process. They differ from each other by how they modify the probabilities, yet in both cases the sampling is parameterized by a hash function that preferably depends on a local context. 
These solutions were improved by~\cite{fernandez2023three}, who proposed a better mark detector, in particular by addressing the issue of repetitive patterns. 

\looseness=-1
We have investigated how to apply these discrete watermarking methods to our audio generation pipeline. For this purpose, we need to encode the audio signal back to tokens in order to identify if the mark is present or not. One issue is that \emph{the codec is not idempotent}: if we generate a waveform from tokens and then re-encode it back into tokens, the re-generated tokens are likely to be different from the ones generated with high probability, even if the audio has not suffered any noise addition. We quantify this problem in \hyperref[tab:idempotence]{Table~\ref{tab:idempotence}}. The semantic token is robust to some extent, while the other quantization indices are increasingly less robust as they depend on the previous quantizer level. One key issue is that the tokens do not resist to a moderate temporal shift. This is especially true for the \mimi codec, which is purposely optimized on a perceptual objective, as opposed to a fidelity reconstruction criterion. 


\begin{table}[t]
\caption{\textbf{Idempotence of tokens}.  Probabilities that quantization indices remain identical after decoding and re-encoding the waveform back to tokens, depending on the residual quantizer level. We consider two optional audio post-processing attacks: audio shifted by a time offset of up to half the sampling period ($\Delta T$=40ms), and re-encoding with RVQGAN. 
All results are averaged over 1000 generated sequences of 1 minute.  
\label{tab:idempotence}
}
\centering 
{%
\footnotesize
\resizebox{\textwidth}{!}{
\begin{tabular}{lccr@{}c@{\ }ccccccc}
\toprule
 & \multicolumn{2}{c}{attacks} & RQ level $\rightarrow$    & $k=1$ & $k=2$ & $k=3$ & $k=4$  & $k=5$ & $k=6$ & $k=7$ & $k=8$ \\   
\cmidrule(l){2-3}
$\downarrow$ codec       & $\Delta T$ & RVQGAN  & & (semantic)  & \multicolumn{6}{c}{}  &\\ 
\midrule
\multirow{4}{*}{Basic}   
   &  0     &  &  &  0.798 & 0.783 & 0.560 & 0.483 & 0.421 & 0.407 & 0.369 & 0.404    \\ 
   &  10ms  &  &  &  0.766 & 0.495 & 0.255 & 0.206 & 0.180 & 0.173 & 0.144 & 0.193  \\ 
   &  20ms  &  &  &  0.682 & 0.390 & 0.220 & 0.180 & 0.158 & 0.154 & 0.129 & 0.172  \\ 
   &  40ms  &  &  &  0.503 & 0.329 & 0.182 & 0.146 & 0.128 & 0.125 & 0.107 & 0.156  \\ 
\midrule
\multirow{8}{*}{\mimi}
   &  0     &  &  & 0.766 & 0.550 & 0.372 & 0.352 & 0.293 & 0.297 & 0.264 & 0.303  \\
   &  10ms  &  &  & 0.731 & 0.376 & 0.206 & 0.176 & 0.152 & 0.154 & 0.132 & 0.182  \\
   &  20ms  &  &  & 0.653 & 0.307 & 0.171 & 0.146 & 0.121 & 0.126 & 0.106 & 0.159  \\
   &  40ms  &  &  & 0.483 & 0.267 & 0.160 & 0.137 & 0.116 & 0.121 & 0.102 & 0.150   \\
\cmidrule(l){2-12}
   &  0    & \checkmark &  & 0.741 & 0.409 & 0.221 & 0.198 & 0.150 & 0.154 & 0.134 & 0.173   \\
   &  10ms & \checkmark &  & 0.702 & 0.281 & 0.148 & 0.133 & 0.118 & 0.117 & 0.100 & 0.136   \\
   &  20ms & \checkmark &  & 0.633 & 0.228 & 0.126 & 0.114 & 0.098 & 0.097 & 0.084 & 0.119   \\
   &  40ms & \checkmark &  & 0.450 & 0.197 & 0.120 & 0.113 & 0.104 & 0.102 & 0.086 & 0.112   \\ 
\bottomrule
\end{tabular}}}
\end{table}


\paragraph{Discussion on generative audio watermarking.} The lack of idempotence is problematic for 
the aforementioned sampling-based watermarking methods, as it affects the reliability of the detector when measuring the sampling bias. Noticeably, in order for these methods to work properly, the n-tuples that gives the context to the hash key must be stable enough during several consecutive tokens. Reducing the context length improves the stability but drastically increases the likelihood of producing degenerated audio sequences, similar to the degeneration problem observed by \citep{holtzman2019curious}. 

While we regard this attempt of employing text-based watermarking as a negative result, hereafter we discuss a few potential ways for  circumventing the aforementioned problem of token stability though re-encoding:
\begin{itemize}
    \item Marking only the RQ first levels improves the stability. In our preliminary experiments, using these indices as context in the hash function, and limiting the dependence on previous timestamps, significantly increases the stability (although not sufficiently). \\[-1.8em]
    \item The idempotence could be improved by adding a specific loss in the discrete latent space, such that the audio tokens are stable through auto-encoding. \\[-1.8em]
    \item Potentially this auto-encoding could be learned to be resilient to signal transformation, similar to what is proposed when learning  image watermarking based on neural networks~\citep{zhu2018hidden,fernandez2022watermarking}. In view of our analysis, adding some tolerance to moderate temporal shift is especially important. \\[-1.8em]
    \item The text could be marked instead of the audio. One downside is that text is a lower-capacity channel for adding a mark, and would not be sufficient for short conversations. Another problem is that detecting the mark requires a reliable transcription. 
\end{itemize}
Last but not least, some exploration is needed to ensure that it is not trivial to remove the watermarking procedure when open-sourcing a model. As an example, the only thing to remove the watermark with the implementation associated with the stable diffusion model was to comment a line of code.\footnote{\url{https://github.com/Stability-AI/stablediffusion/blob/main/scripts/txt2img.py\#L363}} 
A promising work in this direction is the study by \cite{sander2024watermarking}, who show that it is possible to detect when a model has been trained on watermarked text. 
A method exploiting this observation has just been shared by \cite{san2024latent}: the watermarking is implicitly added through the training data, in the spirit of ``radioactive data'' by \cite{sablayrolles2020radioactive}. 
